\documentclass[11pt]{article}
\usepackage[a4paper,margin=24mm]{geometry}
\usepackage{amsmath,amssymb,amsthm,microtype,booktabs}
\usepackage[colorlinks=true,urlcolor=blue,linkcolor=blue]{hyperref}
\newtheorem{theorem}{Theorem}
\newtheorem{lemma}[theorem]{Lemma}
\newtheorem{corollary}[theorem]{Corollary}
\newcommand{\E}{\mathbb E}
\newcommand{\Pp}{\mathbb P}
\newcommand{\Var}{\operatorname{Var}}
\newcommand{\Cov}{\operatorname{Cov}}
\title{Near-linear, but never linear:\\a universal ceiling for consecutive log-concavity failures}
\author{Research draft prepared for Jan Mikul\'ik with ChatGPT}
\date{28 September 2026}
\begin{document}
\maketitle
\begin{abstract}
For every fixed maximum degree $\Delta$, consecutive runs of strict
log-concavity failures in forest independence polynomials have sublinear
length. We prove the effective upper bound
$M_\Delta(N)<6N/\log_2\log_{\Delta-1}N$ for
$N\ge(\Delta-1)^{2^{14}}$, where $\Delta\ge3$ and $M_\Delta(N)$ maximizes
over forests with at most $N$ vertices. The argument combines an exact
positive-coefficient identity converting a convex probability window
into a variance lower bound, entropy control of the appropriate activity,
and covariance decay along tree paths. At any fixed proportional distance
below the independence number, the run length instead satisfies an
explicit polynomial saving. Together with the companion degree-65
construction, this gives exponent one but zero asymptotic density.
\end{abstract}

\section{The universal statement}
Let $i_k(F)$ count independent sets of size $k$ in a finite forest $F$,
and let $\alpha(F)$ be its independence number. A strict failure occurs
at $k$ if $i_k(F)^2<i_{k-1}(F)i_{k+1}(F)$. A run is a consecutive
interval of such indices, and its length is the number of failure
indices, not the number of coefficients in the enlarged window.
For integer $\Delta\ge3$, define
\[
 M_\Delta^{\mathrm{for}}(N)=
 \max_{\substack{|V(F)|\le N\\\Delta(F)\le\Delta}}
 \{\text{length of a strict failure run in }F\}.
\]
The maximum is zero if no failure occurs. Restricting the maximum to
trees gives $M_\Delta^{\mathrm{tree}}(N)\le M_\Delta^{\mathrm{for}}(N)$.

\begin{theorem}\label{thm:finite}
Let $b=\Delta-1\ge2$. For each integer $m\ge1$, set
\begin{equation}\label{eq:thresholds}
 \Lambda_m=4^m,\qquad K_m=\lceil\log_2(12m^2)\rceil,\qquad
 R_m=(\Lambda_m+1)K_m.
\end{equation}
Every $n$-vertex forest of maximum degree at most $\Delta$ with
\[
 n\ge n_{\Delta,m}:=36m^2b^{R_m}
\]
has no failure run of length $L$ satisfying $L+1\ge n/m$.
Consequently, for every
\begin{equation}\label{eq:Nthreshold}
 N\ge N_{\Delta,m}:=36m^3b^{R_m},
 \qquad M_\Delta^{\mathrm{for}}(N)<N/m.
\end{equation}
\end{theorem}

\begin{corollary}\label{cor:global}
For each $\Delta\ge3$,
\begin{equation}\label{eq:global}
 \boxed{M_\Delta^{\mathrm{for}}(N)
 <\frac{6N}{\log_2\log_{\Delta-1}N}
 \quad\text{for }N\ge(\Delta-1)^{2^{14}}.}
\end{equation}
In particular $M_\Delta^{\mathrm{for}}(N)=o(N)$.
For $\Delta=65$, the displayed threshold is the explicit integer
$N\ge2^{98304}$.
\end{corollary}

The large threshold is deliberately conservative. It is a closed
formula, not an unevaluated ``sufficiently large'' constant. Neither
this upper bound nor its constant is claimed to be asymptotically sharp.
All probabilistic estimates used below hold at finite order and finite
positive activity; there is no local limit theorem in this proof.

\section{An exact variance certificate for a convex window}
\begin{lemma}\label{lem:window}
Let $X$ be integer-valued. Suppose its probabilities $p_j$ are convex
on the integer interval $[a,b]$, meaning
$p_{j-1}-2p_j+p_{j+1}\ge0$ for $a<j<b$.
Put $d=b-a\ge2$ and $c=(a+b)/2$. Then
\begin{equation}\label{eq:variance-lower}
 \E(X-c)^2\ge\frac{d(d+2)}{12}.
\end{equation}
Thus if $\E X=c$, the same bound holds for $\Var X$.
\end{lemma}
\begin{proof}
For an arbitrary sequence $u_0,\ldots,u_d$, the exact identity
\begin{equation}\label{eq:dual}
 \begin{split}
 &\sum_{j=0}^{d}\left\{(j-d/2)^2-\frac{d(d+2)}{12}\right\}u_j\\
 &\quad=\sum_{k=1}^{d-1}
 \underbrace{\frac{k(k+1)(d-k)(d-k+1)}{12}}_{C_{d,k}>0}
 (u_{k-1}-2u_k+u_{k+1})
 \end{split}
\end{equation}
follows by collecting coefficients. Specifically, with $C_{d,0}=C_{d,d}=0$,
\[
 C_{d,k-1}-2C_{d,k}+C_{d,k+1}
 =(k-d/2)^2-\frac{d(d+2)}{12}
 \quad(1\le k\le d-1),
\]
and the two endpoint coefficients agree as well. The companion program
verifies these as polynomial identities in two indeterminates, rather
than checking finitely many window lengths.

Apply~\eqref{eq:dual} with $u_j=p_{a+j}$. Its right side is nonnegative.
For an integer $j$ outside $[a,b]$, $|j-c|\ge d/2+1$ and
\[
 (j-c)^2-\frac{d(d+2)}{12}
 \ge\frac{(d+2)(d+3)}6>0.
\]
Adding the outside probabilities proves~\eqref{eq:variance-lower}.
\end{proof}

\paragraph{From a failure run to the lemma.}
Suppose $i_k^2<i_{k-1}i_{k+1}$ for $k=u,\ldots,v$, a run of length
$L=v-u+1$. Set
\begin{equation}\label{eq:window}
 a=u-1,\quad b=v+1,\quad d=L+1,\quad c=(u+v)/2.
\end{equation}
For any $\lambda>0$, the tilted probabilities
\[
 p_k=\Pp_\lambda(X=k)=\frac{i_k\lambda^k}{Z_F(\lambda)},\qquad
 Z_F(\lambda)=\sum_j i_j\lambda^j,
\]
obey the same strict log-convexity inequalities on this window. Since
$p_{k-1}+p_{k+1}\ge2\sqrt{p_{k-1}p_{k+1}}>2p_k$, they are strictly
convex there.
The function $\E_\lambda X$ increases continuously from $0$ to
$\alpha(F)$ as $\lambda$ increases from $0$ to infinity, because
\[
 \frac{d}{d\log\lambda}\E_\lambda X=\Var_\lambda X>0.
\]
There is therefore a unique $\lambda$ for which $\E_\lambda X=c$.
At that activity, Lemma~\ref{lem:window} yields
\begin{equation}\label{eq:run-variance}
 \Var_\lambda X\ge\frac{(L+1)(L+3)}{12}>\frac{(L+1)^2}{12}.
\end{equation}
This argument needs no assumptions on the coefficients outside the run.

\section{Entropy controls the activity}
For the Gibbs distribution on independent sets, its Shannon entropy in
nats is
\[
 H=\log Z_F(\lambda)-(\E_\lambda X)\log\lambda.
\]
Since $Z_F(\lambda)\ge\lambda^{\alpha(F)}$ and the number of independent
sets is at most $2^n$,
\begin{equation}\label{eq:entropy}
 (\alpha(F)-\E_\lambda X)\log\lambda\le H\le n\log2.
\end{equation}
This inequality holds also when $\lambda<1$.
For the centered activity associated with~\eqref{eq:window},
$\alpha(F)-c\ge d/2$, since $b\le\alpha(F)$. Hence
\begin{equation}\label{eq:activity}
 \lambda\le 2^{2n/d}.
\end{equation}
In particular $d\ge n/m$ implies $\lambda\le4^m=\Lambda_m$.

\section{Covariance decay on every finite forest}
\begin{lemma}\label{lem:cov}
Let $F$ have $n$ vertices and maximum degree at most $\Delta\ge3$.
Put $b=\Delta-1$ and $q=\lambda/(1+\lambda)$. For every integer $R\ge0$,
\begin{equation}\label{eq:variance-upper}
 \Var_\lambda X\le\frac{3n b^R+n^2q^R}{4}.
\end{equation}
\end{lemma}
\begin{proof}
Write $S_v=\mathbf1_{\{v\in I\}}$. Root a tree component at $u$.
For a directed parent--child edge $w\to z$, let $q_z$ be the occupation
probability of $z$ conditional on $S_w=0$. The separated subtree has
root odds at most $\lambda$, so $0<q_z\le q$. The exact transition is
\[
 \E(S_z\mid S_w)=q_z(1-S_w).
\]
Removing an edge separates the component, giving the Markov property
along any path. If $u=v_0,v_1,\ldots,v_t=v$, composition of the affine
conditional expectations therefore gives
\begin{equation}\label{eq:covariance}
 \Cov(S_u,S_v)=(-1)^t\Var(S_u)\prod_{j=1}^t q_{v_j},
 \qquad |\Cov(S_u,S_v)|\le\tfrac14q^t.
\end{equation}
Distinct components are independent, so their covariances are zero.

The ball of radius $R$ about a vertex has at most
\[
 1+\Delta\sum_{j=0}^{R-1}(\Delta-1)^j\le3b^R
\]
vertices (the sum is empty when $R=0$). For $R\ge0$, the difference
between $3b^R$ and the geometric expression, multiplied by $b-1$, is
$(2b-4)b^R+2\ge0$.
Sum~\eqref{eq:covariance} over ordered pairs. There are at most
$3nb^R$ near pairs; all other pairs have covariance modulus at most
$q^R/4$, or zero. Bounding the variance by the sum of covariance moduli
proves~\eqref{eq:variance-upper}.
\end{proof}

\section{Proof of the finite threshold and its explicit envelope}
Suppose $d=L+1\ge n/m$ and choose the centered activity. By
\eqref{eq:activity}, $q\le\Lambda_m/(1+\Lambda_m)$. The integer
binomial inequality
\[
 (1+1/\Lambda_m)^{\Lambda_m+1}
 \ge1+(\Lambda_m+1)/\Lambda_m>2
\]
implies, with the exact radius in~\eqref{eq:thresholds},
\begin{equation}\label{eq:qR}
 q^{R_m}<2^{-K_m}\le\frac1{12m^2}.
\end{equation}
If $n\ge36m^2b^{R_m}$, Lemma~\ref{lem:cov} gives
\[
 \Var_\lambda X
 \le\frac{n^2}{48m^2}+\frac{n^2}{48m^2}
 =\frac{n^2}{24m^2}
 <\frac{d^2}{12},
\]
contradicting~\eqref{eq:run-variance}. This proves the exact-order
assertion in Theorem~\ref{thm:finite}.
For $N\ge m n_{\Delta,m}$, a forest with $n<n_{\Delta,m}$ has every
run shorter than $N/m$ simply because $L\le n$. Forests with
$n\ge n_{\Delta,m}$ satisfy $L<n/m\le N/m$. This proves the
at-most-order assertion~\eqref{eq:Nthreshold}. \qed

\paragraph{Deriving the closed $N/\log\log N$ bound.}
For every integer $m\ge1$,
\[
 K_m\le4m,\qquad m\le2^m,\qquad
 R_m\le8m4^m\le2^{3m+3}.
\]
The first inequality follows from $12m^2\le16^m$: it holds at $m=1$,
and $16m^2-(m+1)^2>0$ for $m\ge1$ preserves it. Also
$36m^3\le2^{6+3m}$ and $6+3m\le2^{3m+3}$. Consequently
\begin{equation}\label{eq:threshold-envelope}
 N_{\Delta,m}=36m^3b^{R_m}
 \le b^{R_m+6+3m}\le b^{2^{3m+4}}.
\end{equation}
These inequalities are proved by elementary induction, and their
algebraic witnesses are included in the verifier.

For $N\ge b^{2^{14}}$, let $A=\log_2\log_bN\ge14$ and
$m=\lfloor(A-4)/3\rfloor$. Then $3m+4\le A$ and
\[
 m\ge\frac{A-7}{3}\ge\frac A6.
\]
Equation~\eqref{eq:threshold-envelope} permits applying
Theorem~\ref{thm:finite}, yielding
$M_\Delta^{\mathrm{for}}(N)<N/m\le6N/A$. This proves
Corollary~\ref{cor:global}. \qed

\section{A polynomial barrier away from the final coefficient}
The same argument also forces runs of exponent tending to one to
move towards the independence number.

\begin{theorem}\label{thm:bulk}
Fix $\Delta\ge3$, put $b=\Delta-1$, and let $0<\delta<1$. Define
\[
 \varepsilon_{\Delta,\delta}
 =\frac1{8(1+2^{2/\delta})\log b}>0.
\]
If a failure run of length $L$ in an $n$-vertex forest has midpoint
$c\le(1-\delta)\alpha(F)$ and $n\ge b^4$, then
\begin{equation}\label{eq:bulk}
 \boxed{L+1<\sqrt{12}\,n^{1-\varepsilon_{\Delta,\delta}}.}
\end{equation}
\end{theorem}
\begin{proof}
A forest is bipartite, so $\alpha(F)\ge n/2$. At the centered activity,
\eqref{eq:entropy} gives
$\lambda\le\Lambda:=2^{2/\delta}$.
Take $R=\lfloor\tfrac12\log_b n\rfloor$; since $n\ge b^4$,
\[
 b^R\le\sqrt n,\qquad R\ge\frac{\log n}{4\log b}.
\]
The elementary bound
$\log(1+1/\Lambda)\ge1/(1+\Lambda)$ yields
\[
 q^R\le n^{-\eta},\qquad
 \eta=\frac1{4(1+\Lambda)\log b}=2\varepsilon_{\Delta,\delta}.
\]
Here $0<\eta<1/2$. Thus~\eqref{eq:variance-upper} gives
\[
 \Var_\lambda X\le\tfrac34n^{3/2}+\tfrac14n^{2-\eta}
 \le n^{2-\eta}.
\]
Combine this with~\eqref{eq:run-variance}.
\end{proof}

In particular, for any sequence of bounded-degree forests with
$n\to\infty$ and failure runs of length $L=n^{1-o(1)}$, their
midpoints satisfy $c/\alpha(F)\to1$. Otherwise a fixed $\delta>0$
and an infinite subsequence would contradict~\eqref{eq:bulk}.
By Corollary~\ref{cor:global}, $L=o(n)$; since $\alpha(F)\ge n/2$,
both endpoints of each such run also have ratio tending to one with
$\alpha(F)$. Thus near-linear exponents cannot persist in a fixed
proportional bulk interval, under any bounded-degree construction.

\section{Consequence for the degree-65 construction}
The companion manuscript~\cite{companion} proves
$M_{65}^{\mathrm{tree}}(N)=N^{1-o(1)}$ using even-path seeds and
64-fold branching. Combining that lower bound with the independent
upper bound proved here gives
\begin{equation}\label{eq:two-limits}
 \boxed{
 \lim_{N\to\infty}\frac{\log M_{65}^{\mathrm{tree}}(N)}{\log N}=1,
 \qquad
 \lim_{N\to\infty}\frac{M_{65}^{\mathrm{tree}}(N)}N=0.}
\end{equation}
The first limit uses the companion construction. The second limit,
and the explicit upper bounds in this paper, do not depend on it.
The linear-run target is therefore impossible at any fixed maximum
degree. Equation~\eqref{eq:two-limits} does not determine the exact
subpolynomial factor between the upper and lower bounds.

\section{Verification scope and literature}
The proof above is finite and self-contained. The file
\texttt{verify\_sublinear\_ceiling.py} uses only the Python standard
library and exact rational arithmetic. Running
\begin{verbatim}
python verify_sublinear_ceiling.py
\end{verbatim}
writes \texttt{sublinear\_ceiling\_certificate.json}. The executed run
passed 18 universal algebra/constant checks and an independent exact
enumeration audit of 1364 ordered covariance pairs, with 5245 checks
in total. The algebra checks include the two-variable polynomial
identity~\eqref{eq:dual} and induction witnesses for
\eqref{eq:threshold-envelope}. They are not numerical samples of the
indeterminates. The enumeration audit uses five small forests at four
rational activities and is explicitly a finite sanity check. The
probability, entropy, positivity and induction arguments establishing
the universal quantifiers are the written proofs above; no
proof-assistant formalization is claimed.

For orientation, Yu's September 2026 archive~\cite{Yu26} supplies
unbounded consecutive runs and an $\Omega(N^{1/3})$ lower bound.
Fang, Lu, Nevo, Yao and Zheng~\cite{Fang26} establish eventual
unimodality of forest independence sequences via central-interval
log-concavity. Their proof uses a fixed activity range $[1/4,12]$.
The present upper bound concerns the length of an arbitrary run and
allows the activity to depend on its relative length; it does not
extend their local-limit theorem to all activities. The primitives
used here---Gibbs tilting, entropy and tree conditional independence---are
standard. This literature comparison is not a claim of exhaustive
publication priority. This text is a research draft, not an independently
refereed or formally verified result.

\begin{thebibliography}{9}
\bibitem{companion} \emph{Near-linear runs of strict log-concavity failures
at one fixed maximum degree}. Companion research draft prepared for
Jan Mikul\'ik with ChatGPT, 28 September 2026.
Source file: \texttt{fixed\_degree\_nearlinear.tex}.
\bibitem{Yu26} X. Yu, \emph{Unbounded runs of non-log-concavity in
independence polynomials of trees: reproducibility archive}, version
1.1.0, 3 September 2026.
\url{https://doi.org/10.5281/zenodo.22263331}.
\bibitem{Fang26} E. X. Fang, J. Lu, E. Nevo, Y. Yao, H. Zheng,
\emph{Unimodality of Independence Polynomials for Sufficiently Large
Forests}, arXiv:2609.20961v1, 17 September 2026.
\url{https://arxiv.org/abs/2609.20961}.
\end{thebibliography}
\end{document}
